\documentclass[11pt,a4paper]{article}
\usepackage[utf8]{inputenc}
\usepackage[a4paper,top=2.3cm,bottom=2.6cm,left=2.5cm,right=2.5cm,headsep=0.6cm,footskip=1.1cm]{geometry}
\usepackage{amsmath,amssymb}
\usepackage{enumitem}
\usepackage{fancyhdr}
\usepackage{titlesec}
\usepackage{array}
\usepackage{needspace}
\titlespacing*{\section}{0pt}{1.6em}{0.3em}
\usepackage{hyperref}
\hypersetup{colorlinks=true,linkcolor=black,urlcolor=black,
  pdftitle={FZT 203E Midterm problem pool},pdfauthor={Mehmet Özkan}}

\newcommand{\Z}{\mathbb{Z}}
\newcommand{\R}{\mathbb{R}}
\newcommand{\C}{\mathbb{C}}
\newcommand{\gen}[1]{\langle #1\rangle}
\DeclareMathOperator{\ord}{ord}
\DeclareMathOperator{\sgn}{sgn}
\DeclareMathOperator{\Cl}{Cl}
\DeclareMathOperator{\lcm}{lcm}
\DeclareMathOperator{\tr}{tr}
\DeclareMathOperator{\diag}{diag}
\DeclareMathOperator{\im}{im}
\newcommand{\nsub}{\trianglelefteq}
\DeclareMathOperator{\spn}{span}
\newcommand{\ket}[1]{\lvert #1\rangle}
\newcommand{\up}{{\uparrow}}
\newcommand{\dn}{{\downarrow}}
\newcommand{\MM}{\mathsf{M}}

% Solutions are kept out of this public file. The instructor copy defines
% \withsolutions and loads GT_midterm_pool_solutions.tex (not published).
\makeatletter
\newcommand{\defsolution}[2]{\@namedef{sol@#1}{#2}}
\newcommand{\showsolution}[1]{%
  \ifdefined\withsolutions\@ifundefined{sol@#1}{}{%
    \par\vspace{0.6em}\noindent\textbf{Solution.}\par\vspace{0.1em}{\sloppy\@nameuse{sol@#1}\par}\vspace{0.4em}\hrule}\fi}
\makeatother

\pagestyle{fancy}
\fancyhf{}
\renewcommand{\headrulewidth}{0pt}
\fancyfoot[L]{\small FZT 203E $\cdot$ Midterm problem pool\ifdefined\withsolutions\ $\cdot$ instructor copy with solutions, not for distribution\fi}
\fancyfoot[R]{\small\thepage}

\setlength{\parindent}{0pt}
\setlength{\parskip}{0.55em}

\newcounter{poolprob}
\newcommand{\poolproblem}[2]{\par\vspace{1.1em}\refstepcounter{poolprob}\noindent{\bfseries \thepoolprob. #1}%
  \if\relax\detokenize{#2}\relax\else\hfill{\small\textsf{#2}}\fi\par\vspace{0.1em}}
\setlist[enumerate,1]{label=\textbf{(\alph*)},leftmargin=2.2em,itemsep=0.35em,topsep=0.2em}
\setlist[itemize]{label=$\bullet$,itemsep=0.25em,topsep=0.2em}

\begin{document}
\ifdefined\withsolutions\input{GT_midterm_pool_solutions}\fi

{\small FZT 203E $\cdot$ Theory of Groups in Physics $\cdot$ Fall 2026 \hfill Istanbul Technical University\\
Instructor: Prof.\ Mehmet Özkan}

\vspace{0.9em}
{\LARGE\bfseries Midterm problem pool}\ifdefined\withsolutions\hfill{\large\bfseries Instructor copy}\fi

\vspace{0.5em}
The midterm is a three-hour exam of five problems, each made of selected parts of one of the 35 problems below. The problems cover Lectures 1--10 (Weeks 1--5). Problems 1--7 are Problem Set 1, Problems 10--14 are Problem Set 2, and Problems 17 and 24--26 are Problem Set 3.

Justify every answer; a correct answer without reasoning earns partial or no credit.

\section*{\large Conventions}
These are the conventions of the lectures, the problem sets and the course simulations.
\begin{itemize}
  \item The vertices of a regular $n$-gon are labelled $1, 2, \dots, n$ counterclockwise. $r$ is the counterclockwise rotation by $2\pi/n$, and $s$ is the reflection in the axis through vertex 1 and the centre. $D_n$ has the relations $r^n = s^2 = e$ and $srs = r^{-1}$, so $s r^k = r^{-k} s$, and every element has a unique normal form $r^k s^m$ with $0 \le k < n$ and $m \in \{0, 1\}$.
  \item Products are read right to left: in $gh$, the element $h$ acts first. A symmetry is written as a permutation of positions: $i \to j$ when the object at position $i$ is carried to position $j$. Permutations compose the same way: $(12)(23)$ means first $(23)$, then $(12)$. For the triangle, $r = (123)$, $s = (23)$ and $rs = (12)$.
  \item $a$ and $b$ are conjugate if $b = g^{-1} a g$ for some $g \in G$, and $\Cl(a) = \{g^{-1} a g : g \in G\}$ is the class of $a$. A normal subgroup is written $N \nsub G$, and $G/N$ is the quotient group of the cosets $gN$. $\sgn(p) = \pm 1$ is the sign of $p \in S_n$, and $A_n = \ker(\sgn)$.
  \item $\Z_n = \{0, 1, \dots, n-1\}$ under addition mod $n$. $D_2 \cong \Z_2 \times \Z_2$ is the four-element group of half-turns of a box.
  \item Vectors are columns. A representation is a homomorphism $D\colon G \to GL(d, \C)$, $D(gh) = D(g)D(h)$; representations are finite dimensional over $\C$ unless a real space is specified. $D(g)_{ij}$ is the entry in row $i$ and column $j$.
  \item For a finite group, $N = |G|$; the conjugacy classes are $C_1, \dots, C_k$, with $n_i = |C_i|$ elements. The character of $D$ is $\chi_D(g) = \tr D(g)$. For functions on $G$,
  \[
    \langle f, h\rangle_G = \frac{1}{N}\sum_{g \in G} f(g)^* h(g) = \frac{1}{N}\sum_{i=1}^{k} n_i\, f(C_i)^* h(C_i),
  \]
  the second form for class functions. The multiplicity of the irreducible representation $D^{(\alpha)}$ in $D$ is $m_\alpha = \langle \chi_\alpha, \chi_D\rangle_G$.
  \needspace{16\baselineskip}
  \item The plane representation of $D_3 \cong S_3$, with vertex 1 at $(1, 0)$, is
  \[
    D_E(r) = R = \begin{pmatrix} -\tfrac12 & -\tfrac{\sqrt3}{2} \\[2pt] \tfrac{\sqrt3}{2} & -\tfrac12 \end{pmatrix}, \qquad
    D_E(s) = F = \begin{pmatrix} 1 & 0 \\ 0 & -1 \end{pmatrix},
  \]
  using $\cos 120^\circ = -\tfrac12$ and $\sin 120^\circ = \tfrac{\sqrt3}{2}$. Whenever a problem uses a specific angle, the sines and cosines it needs are given.
  The classes of $S_3$ are $\{e\}$, the two 3-cycles $\{r, r^2\}$ (rotations) and the three transpositions $\{s, rs, r^2 s\}$ (reflections). Its irreducible representations are the trivial representation $1$, the sign representation $\sgn$ and $E$, with the characters
  \[
    \begin{array}{c|ccc}
      & \{e\} & \{r, r^2\} & \{s, rs, r^2 s\} \\ \hline
      1 & 1 & 1 & 1 \\
      \sgn & 1 & 1 & -1 \\
      E & 2 & -1 & 0
    \end{array}
  \]
  \item You may use the results of the lectures, among them Lagrange's theorem, the first isomorphism theorem, the finite-group unitarity theorem, both forms of Schur's lemma, the Great Orthogonality Theorem, the orthonormality of irreducible characters, $\sum_\alpha d_\alpha^2 = N$, and that the number of inequivalent irreducible representations equals the number of conjugacy classes.
\end{itemize}

%======================================================================
\section*{\large Part A. Groups, subgroups, cosets and homomorphisms (Lectures 1--4)}

\poolproblem{Group or not?}{Problem Set 1, Problem 1}
\begin{enumerate}
  \item Decide which of the following are groups. For each one that is not, name an axiom that fails and give an explicit counterexample. For each one that is, say whether it is Abelian.
  \begin{enumerate}[label=\arabic*.,leftmargin=1.8em]
    \item $\Z$ with $a * b = a - b$.
    \item The odd integers under addition.
    \item The nonzero integers under multiplication.
    \item $\{1, i, -1, -i\}$ under multiplication of complex numbers.
    \item The real $2 \times 2$ matrices with determinant 1, under matrix multiplication.
    \item $\R \setminus \{-1\}$ with $a * b = a + b + ab$. (Hint: compute $1 + a * b$.)
  \end{enumerate}
  \item Suppose every element of a group $G$ satisfies $g^2 = e$. Show that $G$ is Abelian.
  \item Let $G$ be a finite group. Show that the elements with $g \ne g^{-1}$ come in pairs, so there is an even number of them. Deduce that every group of even order contains an element of order 2, and check this in $D_3$ and $\Z_6$.
\end{enumerate}
\showsolution{1}

\poolproblem{Working in $D_4$}{Problem Set 1, Problem 2}
Take the square with vertices $1, 2, 3, 4$ counterclockwise. Check that, with our conventions, $r = (1234)$ and $s = (24)$.
\begin{enumerate}
  \item Write all eight elements of $D_4$ in normal form and as permutations in cycle notation. For each reflection, describe its axis.
  \item Show that $s r^k = r^{-k} s$ for every integer $k$. Use this to bring $x = s r^3 s r^2 s r$ to normal form, then check your answer by multiplying the permutations.
  \item Find the order of every element of $D_4$.
  \item $D_4$ has exactly ten subgroups. Find them all, say which are cyclic and which are isomorphic to $D_2$, and draw the subgroup lattice: each subgroup joined by a line to the subgroups directly above it that contain it. Explain why there are no others.
\end{enumerate}
\showsolution{2}

\poolproblem{The cyclic group $\Z_{12}$}{Problem Set 1, Problem 3}
\begin{enumerate}
  \item Find the order of every element of $\Z_{12}$. Guess a formula for the order of $k$ in terms of $k$ and 12, and check it on all twelve elements.
  \item Which elements generate $\Z_{12}$?
  \item Find all subgroups of $\Z_{12}$ and draw the subgroup lattice. You may use without proof that every subgroup of a cyclic group is cyclic.
  \item For each subgroup $H$, list the elements $k$ with $\gen{k} = H$. What do the elements that generate the same subgroup have in common?
\end{enumerate}
\showsolution{3}

\poolproblem{Cosets and Lagrange's theorem}{Problem Set 1, Problem 4}
\begin{enumerate}
  \item In $D_4$ let $H = \gen{s} = \{e, s\}$. List the left cosets $gH$ and the right cosets $Hg$. Which left cosets are also right cosets?
  \item Repeat (a) for $K = \gen{r^2}$ and for $R = \gen{r}$. What property of $r^2$ explains your answer for $K$?
  \item A group $G$ has order 20. Which orders are possible for its subgroups, and which for its elements?
  \item A group $G$ has order 15, and an element $g \in G$ satisfies $g^3 \ne e$ and $g^5 \ne e$. Show that $G$ is cyclic.
  \item $H$ and $K$ are subgroups of a group $G$ with $|H| = 4$ and $|K| = 9$. Show that $H \cap K = \{e\}$.
  \item By Lagrange's theorem the order of every element divides $|G|$. Give an example of a group $G$ and a divisor $d$ of $|G|$ such that $G$ has no element of order $d$.
\end{enumerate}
\showsolution{4}

\poolproblem{Homomorphisms}{Problem Set 1, Problem 5}
\begin{enumerate}
  \item Each element $g$ of $D_n$ acts on the plane as a $2 \times 2$ matrix $M(g)$, with $M(gh) = M(g)M(h)$. With vertex 1 on the positive $x$-axis,
  \[
    M(r) = \begin{pmatrix} \cos\frac{2\pi}{n} & -\sin\frac{2\pi}{n} \\[2pt] \sin\frac{2\pi}{n} & \cos\frac{2\pi}{n} \end{pmatrix}, \qquad
    M(s) = \begin{pmatrix} 1 & 0 \\ 0 & -1 \end{pmatrix}.
  \]
  Define $\varphi(g) = \det M(g)$. Show that $\varphi$ is a homomorphism from $D_n$ to the group $\{1, -1\}$ under multiplication, compute $\varphi(r^k s^m)$, and find the kernel and the image of $\varphi$.
  \item Let $U(1) = \{z \in \C : |z| = 1\}$ under multiplication, and let $\psi(z) = z^2$. Show that $\psi$ is a homomorphism from $U(1)$ onto $U(1)$, find its kernel, and show that every element of $U(1)$ has exactly two preimages. More generally, show that the kernel of $z \mapsto z^n$ is isomorphic to $\Z_n$. (In Week 9 the map $SU(2) \to SO(3)$ will be two-to-one in the same way.)
  \item Find all homomorphisms $\Z_4 \to \Z_6$ and all homomorphisms $\Z_6 \to \Z_4$, and give the kernel and image of each. A homomorphism from a cyclic group is fixed by where it sends a generator; check that each candidate is well defined.
  \item For every homomorphism $\varphi\colon G \to G'$ you found in (a) and (c), check that $|\ker\varphi| \cdot |\im\varphi| = |G|$. We will see why this always holds when we meet quotient groups.
\end{enumerate}
\showsolution{5}

\poolproblem{Isomorphic or not?}{Problem Set 1, Problem 6}
\begin{enumerate}
  \item For each pair, either give an explicit isomorphism or prove that none exists. To rule one out, find a property that isomorphisms preserve, such as being Abelian or the orders of elements, that one group has and the other lacks. You may assume that each set below is a group.
  \begin{enumerate}[label=\arabic*.,leftmargin=1.8em]
    \item $\Z_6$ and $S_3$.
    \item $\Z_4$ and $\{1, 2, 3, 4\}$ under multiplication mod 5.
    \item $\Z_4$ and $\{1, 3, 5, 7\}$ under multiplication mod 8.
    \item $(\R, +)$ and $(\R_+, \cdot\,)$, the positive reals under multiplication.
    \item $(\R, +)$ and $(\R \setminus \{0\}, \cdot\,)$.
  \end{enumerate}
  \item Label the elements $1, 3, 5, 7$ of the group in item 3 as positions $1, 2, 3, 4$. Following Cayley's theorem, write multiplication by each element as a permutation in $S_4$. What do the four permutations have in common?
\end{enumerate}
\showsolution{6}

\poolproblem{Adding velocities in special relativity}{Problem Set 1, Problem 7}
In units with $c = 1$, two collinear velocities $u$ and $v$ combine to
\[
  u \oplus v = \frac{u + v}{1 + uv}.
\]
Let $G = (-1, 1)$, the possible velocities of a massive particle.
\begin{enumerate}
  \item Show that $u \oplus v \in G$ whenever $u, v \in G$. (Hint: factor $1 - (u \oplus v)$ and $1 + (u \oplus v)$.)
  \item Find the identity and the inverse of $u$. Is $\oplus$ commutative?
  \item The rapidity $\eta$ of a velocity $u$ is defined by $u = \tanh\eta$. Use
  \[
    \tanh(\alpha + \beta) = \frac{\tanh\alpha + \tanh\beta}{1 + \tanh\alpha\tanh\beta}
  \]
  to show that $T(\eta) = \tanh\eta$ satisfies $T(\alpha + \beta) = T(\alpha) \oplus T(\beta)$. Since $T$ is a bijection from $\R$ onto $(-1, 1)$, write $u = T(\alpha)$, $v = T(\beta)$, $w = T(\gamma)$ and deduce that $\oplus$ is associative without expanding $(u \oplus v) \oplus w$. Conclude that $(G, \oplus)$ is a group isomorphic to $(\R, +)$.
  \item Compute $0.5 \oplus 0.5$, then check your answer by adding rapidities. (Hint: $\tanh^{-1} u = \tfrac12 \ln\bigl[(1 + u)/(1 - u)\bigr]$.)
  \item With $c$ restored, $u \oplus v = (u + v)/(1 + uv/c^2)$. Which group do you recover as $c \to \infty$?
\end{enumerate}
\showsolution{7}

\poolproblem{Two mirrors: kaleidoscopes and dihedral groups}{}
Identify the plane with $\C$ through $(x, y) \leftrightarrow z = x + iy$, and write $\bar z$ for the complex conjugate of $z$. For real angles $\alpha$ and $\theta$ let
\[
  \sigma_\alpha(z) = e^{2i\alpha}\,\bar z, \qquad \rho_\theta(z) = e^{i\theta} z.
\]
\begin{enumerate}
  \item Show that $\sigma_\alpha$ fixes every point $z = t e^{i\alpha}$ ($t$ real) of the line through 0 at the angle $\alpha$ to the positive $x$-axis and reverses the perpendicular direction, $\sigma_\alpha(i e^{i\alpha}) = -i e^{i\alpha}$: it is the reflection in that line. Show that $\sigma_\alpha^2 = e$, and that $\rho_\theta$ is the counterclockwise rotation by $\theta$.
  \item Show that $\sigma_\alpha \sigma_\beta = \rho_{2(\alpha - \beta)}$, where $\sigma_\beta$ acts first, and that $\rho_\theta \sigma_\alpha = \sigma_{\alpha + \theta/2}$ and $\sigma_\alpha \rho_\theta = \sigma_{\alpha - \theta/2}$. Thus two reflections compose to a rotation by twice the angle between their mirrors. Does the order of the two reflections matter?
  \item A kaleidoscope has two mirrors through 0 at the angle $\pi/n$, $n \ge 2$: $s = \sigma_0$ and $t = \sigma_{\pi/n}$. Show that $r = ts$ is the counterclockwise rotation by $2\pi/n$, that $t = rs$ and that $srs = r^{-1}$. Deduce that $s$ and $t$ generate the $2n$ distinct transformations $r^k s^m$, $0 \le k < n$, $m \in \{0, 1\}$: the group $D_n$, with vertex 1 of the $n$-gon at $z = 1$. Find the mirror line of each reflection $r^k s$.
  \item An object sits at a point $z_0$ strictly inside the wedge $0 < \arg z < \pi/n$ between the mirrors, and the kaleidoscope shows it at the points $g(z_0)$, $g \in D_n$. Show that these $2n$ points are all different (hint: if $g(z_0) = h(z_0)$, then $h^{-1}g$ fixes $z_0$), so that the viewer sees the object and $2n - 1$ images. How many images do mirrors at $60^\circ$ and at $45^\circ$ produce?
  \item Now let the angle between the mirrors be any $\alpha$ with $0 < \alpha < \pi$, with $s = \sigma_0$ and $t = \sigma_\alpha$. Show that the group generated by $s$ and $t$ is finite if and only if $\alpha/\pi$ is rational. For $\alpha = 2\pi/5$, find the order of this group and show that it consists of the same transformations as the group for $\alpha = \pi/5$.
\end{enumerate}
\showsolution{8}

\poolproblem{Groups of order six: a fan and the ammonia molecule}{}
The rotations of a fan with six identical, twisted blades form a cyclic group of order 6; no reflection is a symmetry because the blades are twisted. The symmetries of an ammonia molecule NH$_3$, the rotations by $0$ and $\pm 120^\circ$ about the axis through the nitrogen atom and the reflections in the three planes that contain this axis and one hydrogen atom, form a group isomorphic to $S_3$. This problem shows that there are no other groups of order six. Let $G$ be a group with $|G| = 6$.
\begin{enumerate}
  \item Which orders can the elements of $G$ have? Show that if $G$ has an element of order 6, then $G \cong \Z_6$.
  \item From now on assume that $G$ has no element of order 6. Show that $G$ has an element $a$ of order 3. (Hint: otherwise $g^2 = e$ for every $g \in G$. Show that $G$ is then Abelian, and that two distinct elements $a, b \ne e$ then give a subgroup $\{e, a, b, ab\}$ of order 4.)
  \item Let $H = \gen{a}$ and choose $b \in G$ with $b \notin H$. Show that the right cosets of $H$ are $H$ and $Hb$, so that $G = \{e, a, a^2, b, ab, a^2 b\}$.
  \item Show that $b^2 = e$. (Hint: show that $b^2$ cannot lie in $Hb$. If $b^2 = a$ or $b^2 = a^2$, what is the order of $b$?)
  \item Show that $bab^{-1} \in H$, and then that $bab^{-1} = a^{-1}$. (Hint: if $ab = ba$, find the order of $ab$.)
  \item Conclude that $a^k b^m \mapsto r^k s^m$ is an isomorphism $G \to D_3 \cong S_3$, so that every group of order six is isomorphic to $\Z_6$ or to $S_3$. In the ammonia molecule, which symmetries can play the roles of $a$ and $b$?
\end{enumerate}
\showsolution{9}

%======================================================================
\section*{\large Part B. Direct products, classes, quotients and permutations (Lectures 4--6)}

\poolproblem{Direct products, lattices, parity and time reversal}{Problem Set 2, Problem 1}
\begin{enumerate}
  \item Let $g \in G$ and $k \in K$ have finite orders. Show that the order of $(g, k)$ in $G \times K$ is $\lcm\bigl(\ord(g), \ord(k)\bigr)$. You may use that $g^m = e$ exactly when $\ord(g)$ divides $m$.
  \item Deduce that $\Z_m \times \Z_n \cong \Z_{mn}$ when $\gcd(m, n) = 1$, and that $\Z_m \times \Z_n$ is not cyclic when $\gcd(m, n) > 1$. Here $\gcd(m, n)$ is the greatest common divisor of $m$ and $n$, the largest positive integer that divides both. This generalizes $\Z_2 \times \Z_3 \cong \Z_6$ from Lecture 4.
  \item A particle hops on a $6 \times 9$ lattice with periodic boundary conditions in both directions, so its position $(i, j)$ is an element of $\Z_6 \times \Z_9$. At every time step it is translated by the same $(a, b) \in \Z_6 \times \Z_9$. Starting from $(0, 0)$, how many different sites does it visit for $(a, b) = (1, 1)$, $(2, 3)$ and $(3, 1)$? What is the largest possible number, and can the particle ever visit all 54 sites? Answer the last question also for a $4 \times 9$ lattice.
  \item A particle on a line obeys $m\ddot x = F(x, \dot x)$. Acting on trajectories $t \mapsto x(t)$, define parity and time reversal by
  \[
    (Px)(t) = -x(t), \qquad (Tx)(t) = x(-t).
  \]
  Show that $P^2 = T^2 = e$ and $PT = TP$, so that $\{e, P, T, PT\}$ is a group isomorphic to $\Z_2 \times \Z_2$, not to $\Z_4$.
  \item Call one of these four transformations a symmetry if it maps every solution of the equation of motion to a solution. Explain why the symmetries among $e, P, T, PT$ form a subgroup. Find this subgroup for a spring, $F = -kx$; a spring with friction, $F = -kx - \gamma\dot x$; a falling ball, $F = -mg$; and a falling ball with air drag, $F = -mg - \gamma\dot x$ (all constants positive). Which physical feature destroys $P$, and which destroys $T$? (Hint: first find the condition on $F$ for each of $P$, $T$ and $PT$ to be a symmetry.)
\end{enumerate}
\showsolution{10}

\poolproblem{The symmetries of benzene}{Problem Set 2, Problem 2}
The six carbon atoms of a benzene molecule sit at the vertices of a regular hexagon, labelled $1, \dots, 6$ counterclockwise. Its symmetries in the plane of the ring form $D_6$, with $r = (123456)$ and $s = (26)(35)$. (In Bünemann's realization of $D_n$ by rotations of a prism, the reflections become half-turns; nothing below changes.)
\begin{enumerate}
  \item Show that $r^{-j}(r^k s)\, r^j = r^{k - 2j} s$ and $s^{-1}(r^k s)\, s = r^{-k} s$. Find all conjugacy classes of $D_6$, and check that their sizes add up to 12.
  \item Write each of the six reflections as a permutation, and say whether its mirror passes through two atoms or through the midpoints of two C--C bonds. How do the two classes of reflections differ geometrically? (This distinction returns with the point groups of Week 7.)
  \item Let $Z = \{e, r^3\}$, $N = \{e, r^2, r^4\}$ and $H = \{e, r^2, r^4, s, r^2 s, r^4 s\}$. Show that $Z$, $N$ and $H$ are normal subgroups of $D_6$ and that $\{e, s\}$ is not. Show that $r^3$ commutes with every element of $D_6$. What does $r^3$ do to a point of the plane?
  \item Show that $H$ consists of the symmetries of the hexagon that map the triangle of atoms $1, 3, 5$ to itself, and that $H \cong D_3$. Show that every $g \in D_6$ can be written in exactly one way as $g = hz$ with $h \in H$ and $z \in Z$, and that $(h, z) \mapsto hz$ is an isomorphism $H \times Z \to D_6$. Conclude that $D_6 \cong D_3 \times \Z_2$ and that $D_6 / Z \cong D_3$.
  \item List the cosets of $N$, compute the square of each coset, and decide which group of order four $D_6/N$ is. Two symmetries lie in the same coset of $N$ exactly when they give the same answers to two yes-or-no questions. What are the questions?
\end{enumerate}
\showsolution{11}

\poolproblem{Angles, spokes and the wagon-wheel effect}{Problem Set 2, Problem 3}
Since $(\R, +)$ is Abelian, every $c\Z = \{ck : k \in \Z\}$, $c > 0$, is a normal subgroup and $\R/c\Z$ is a group: the continuous counterpart of $\Z/N\Z \cong \Z_N$ from Lecture 5.
\begin{enumerate}
  \item Show that $E(\theta) = e^{i\theta}$ is a homomorphism from $(\R, +)$ onto $U(1) = \{z \in \C : |z| = 1\}$ with kernel $2\pi\Z$, and conclude that $\R/2\pi\Z \cong U(1)$. Show that every coset $\theta + 2\pi\Z$ contains exactly one number in $(-\pi, \pi]$. This is the precise sense in which an angle is defined modulo $2\pi$.
  \item A wheel has $n$ identical, equally spaced spokes, so turning it by $\theta$ and by $\theta + 2\pi/n$ gives the same picture. Show that $\theta \mapsto e^{in\theta}$ is a homomorphism from $\R$ onto $U(1)$ with kernel $(2\pi/n)\Z$. Thus the pictures of the wheel are labelled by $\R/(2\pi/n)\Z \cong U(1)$, and each coset has exactly one representative in $(-\pi/n, \pi/n]$.
  \item A film camera takes 24 frames per second, so a wheel turning at angular velocity $\omega$ turns by $\omega/24$ between frames, and the film records only the coset $\omega/24 + (2\pi/n)\Z$. Assume that at each frame the audience sees the spokes turn by the representative of this coset in $(-\pi/n, \pi/n]$, the smallest turn consistent with the pictures. A stagecoach wheel of radius $0.5$\,m with 12 spokes rolls without slipping, so $\omega = v/(0.5\,\text{m})$ at speed $v$.
  Up to what speed does the film show the true rotation? At which speeds does the wheel seem to stand still? At $v = 5$\,m/s and at $v = 7$\,m/s, how fast and in which direction does the wheel seem to turn?
  \item The same quotient labels waves on a crystal. The atoms of a one-dimensional crystal sit at $x_j = ja$, $j \in \Z$, and a wave with wave number $k$ has the phase $e^{ikx_j}$ at atom $j$. Show that wave numbers $k$ and $k'$ give the same phase at every atom if and only if $k - k' \in (2\pi/a)\Z$. Lattice waves are therefore labelled by $\R/(2\pi/a)\Z$, with one representative of each coset in $(-\pi/a, \pi/a]$, the first Brillouin zone of solid-state physics. Which wave number in this interval describes the same lattice wave as $k = 7\pi/(2a)$?
\end{enumerate}
\showsolution{12}

\poolproblem{Shuffles, recurrence and the 15 puzzle}{Problem Set 2, Problem 4}
\begin{enumerate}
  \item A deck of 8 cards fills positions $1, \dots, 8$ from the top. A perfect out-shuffle cuts the deck into halves, positions 1--4 and 5--8, and interleaves them with the top card staying on top: the new deck holds, from the top, the cards formerly at positions $1, 5, 2, 6, 3, 7, 4, 8$. An in-shuffle interleaves the other way: $5, 1, 6, 2, 7, 3, 8, 4$. Write each shuffle as a permutation of positions in disjoint-cycle notation. How many out-shuffles, and how many in-shuffles, restore the original order? Find the sign of each shuffle.
  \item $N$ labelled particles occupy $N$ sites, and a deterministic, reversible rule moves them in the same way at every time step, so one step is a permutation $p \in S_N$. Explain why the starting configuration recurs, and why it first recurs after $\ord(p)$ steps: a finite version of Poincaré recurrence in mechanics. When does a single particle first return to its starting site? Show that the recurrence time can exceed $N$: find the largest order of an element of $S_7$, and write down an element of $S_{10}$ of order 30.
  \item In the 15 puzzle, fifteen tiles numbered 1--15 and one empty square fill a $4 \times 4$ frame; a move slides a tile next to the empty square into it. Number the squares 1--16 row by row from the top left, and treat the empty square as a sixteenth tile. Start with tile $i$ on square $i$, so the empty square is square 16, in the bottom right corner. Show that each move is a transposition. By colouring the squares like a chessboard, show that the empty square can return to square 16 only after an even number of moves. Deduce that every arrangement reachable with the empty square back at square 16 is an even permutation of the tiles $1, \dots, 15$.
  \item Can the puzzle reach the arrangement in which only tiles 14 and 15 are exchanged? Can it reach the arrangement in which the tiles read $15, 14, \dots, 1$ along the rows, with the empty square still at square 16?
\end{enumerate}
\showsolution{13}

\poolproblem{Rotations of a cube}{Problem Set 2, Problem 5}
Place a cube with its vertices at $(\pm1, \pm1, \pm1)$, so that the coordinate axes pass through the centres of opposite faces. Let $G$ be the group of rotations that map the cube to itself; reflections are not included. $G$ is the rotational symmetry group of a cubic crystal such as rock salt (in the Schoenflies notation of Week 7 it is called $O$). A rotation is given by its effect on points: for example, $(x, y, z) \mapsto (-y, x, z)$ is the counterclockwise quarter-turn about the $z$-axis. The four body diagonals join opposite vertices: $d_1$ passes through $(1, 1, 1)$, $d_2$ through $(-1, 1, 1)$, $d_3$ through $(-1, -1, 1)$ and $d_4$ through $(1, -1, 1)$.
\begin{enumerate}
  \item Every rotation axis of the cube passes through the centres of two opposite faces, through two opposite vertices, or through the midpoints of two opposite edges. For each kind, give the number of axes and the angles of the rotations about each one that map the cube to itself. Check that, together with $e$, there are 24 rotations. Why can there be no others? (Hint: a rotation is fixed once you know which face it moves to the top and how that face is turned.)
  \item Each $g \in G$ permutes the four diagonals. Let $\varphi(g) \in S_4$ be the permutation with $i \to j$ when $g$ carries $d_i$ to $d_j$. Explain why $\varphi(gh) = \varphi(g)\varphi(h)$. Find $\varphi(g)$ for
  \[
    R\colon (x, y, z) \mapsto (-y, x, z), \qquad
    C\colon (x, y, z) \mapsto (z, x, y), \qquad
    B\colon (x, y, z) \mapsto (y, x, -z).
  \]
  $R$ is the quarter-turn about the $z$-axis, $C$ a turn by $120^\circ$ about $d_1$, and $B$ the half-turn about the line through $(1, 1, 0)$ and $(-1, -1, 0)$.
  \item Show that $\ker\varphi = \{e\}$. (Hint: if $g$ maps every diagonal to itself, it sends each vertex $v$ to $v$ or to $-v$. Compare the distance between two neighbouring vertices before and after. Is $(x, y, z) \mapsto (-x, -y, -z)$ a rotation?) Conclude that $\varphi$ is an isomorphism, $G \cong S_4$.
  \item An isomorphism preserves the order of every element. Use this and the identity $\varphi(R^2) = \varphi(R)^2$ to match each kind of rotation in (a) with a conjugacy class of $S_4$ from the table of Lecture 6, and check that the numbers agree. Conclude that two rotations are conjugate in $G$ exactly when they are of the same kind. In particular, the quarter-turns by $+90^\circ$ and $-90^\circ$ about the $z$-axis are conjugate: find $g \in G$ with $g^{-1} R g = R^{-1}$, and explain the result physically.
  \item In a methane molecule CH$_4$, the carbon atom sits at the centre of the cube and the four hydrogen atoms at the alternate vertices $(1, 1, 1)$, $(1, -1, -1)$, $(-1, 1, -1)$ and $(-1, -1, 1)$, which form the set $T_+$; the other four vertices form $T_-$. Show that every $g \in G$ maps $T_+$ either onto itself or onto $T_-$ (hint: any two points of $T_+$ are joined by a diagonal of a face), and that setting $\chi(g) = +1$ or $-1$ accordingly defines a homomorphism $\chi\colon G \to \{+1, -1\}$. Evaluate $\chi$ on $R$, $R^2$, $C$ and $B$, and use (d) to show that $\chi(g) = \sgn\varphi(g)$ for every $g \in G$. Conclude that the rotations of the cube that map the methane molecule to itself form a group isomorphic to $A_4$. What does the sign of $\varphi(g)$ mean geometrically?
  \item Each $g \in G$ also permutes the three coordinate axes, which defines a homomorphism $\psi\colon G \to S_3$. Find $\psi(R)$ and $\psi(C)$, and show that $\psi$ is onto. Use the first isomorphism theorem to find $|\ker\psi|$, and identify $\ker\psi$ without further search. Deduce that $V = \{e, (12)(34), (13)(24), (14)(23)\}$ is a normal subgroup of $S_4$ and that $S_4/V \cong S_3$. Which group from Week 1 is $\ker\psi$?
\end{enumerate}
\showsolution{14}

\poolproblem{Bosons and fermions: one-dimensional exchange representations}{}
A state of $N \ge 2$ identical particles is described by a wavefunction $\Psi(x_1, \dots, x_N) \ne 0$, where $x_i$ is the position of particle $i$ together with its spin (and any other internal label it carries). For $p \in S_N$ define the relabelled wavefunction
\[
  \Psi_p(x_1, \dots, x_N) = \Psi\bigl(x_{p(1)}, \dots, x_{p(N)}\bigr).
\]
Because the particles are identical, no measurement distinguishes $\Psi_p$ from $\Psi$. We assume, as is standard, that $\Psi_p$ then describes the same physical state as $\Psi$, so that $\Psi_p = c(p)\,\Psi$ for a nonzero complex number $c(p)$. This is an extra assumption about relabelling, not a consequence of identity alone: it allows only one-dimensional representations of $S_N$, as (a) shows.
\begin{enumerate}
  \item Show that $(\Psi_p)_q = \Psi_{qp}$, and deduce that $c(qp) = c(q)\,c(p)$: $c$ is a homomorphism from $S_N$ to the group $\C^\times = \C \setminus \{0\}$ under multiplication.
  \item Show that $c(\tau)^2 = 1$ for every transposition $\tau$, and that any two transpositions are conjugate in $S_N$. Deduce that $c$ takes the same value, $+1$ or $-1$, on every transposition.
  \item Every permutation is a product of transpositions. Conclude that either $c(p) = 1$ for every $p$ (bosons) or $c(p) = \sgn(p)$ for every $p$ (fermions). Find $\ker c$ in each case, and show that $S_N/\ker c \cong \Z_2$ for fermions.
  \item For two fermions in the single-particle states $\phi_a$ and $\phi_b$, $\Psi(x_1, x_2) = \phi_a(x_1)\phi_b(x_2) - \phi_a(x_2)\phi_b(x_1)$. Check that $\Psi_{(12)} = -\Psi$, and show that $\Psi = 0$ when $a = b$: two identical fermions cannot occupy the same state (the Pauli exclusion principle).
  \item The same argument classifies the homomorphisms $c\colon D_n \to \C^\times$. Apply $c$ to $srs = r^{-1}$ to show that $c(r) = \pm1$; show also that $c(s) = \pm 1$. Show that $c(r) = 1$ when $n$ is odd, so that $D_n$ has only two such homomorphisms, while for even $n$ all four sign choices give homomorphisms. Explain the difference with the conjugacy classes of the reflections in $D_n$.
\end{enumerate}
\showsolution{15}

\poolproblem{Domain walls on an Ising ring}{}
$N \ge 3$ spins sit on a ring: spin $i$ takes the value $\sigma_i = \pm 1$, and spin $N$ is next to spin 1, so $\sigma_{N+1} = \sigma_1$. A configuration is $\sigma = (\sigma_1, \dots, \sigma_N)$, and configurations are multiplied entry by entry, $(\sigma\tau)_i = \sigma_i \tau_i$. Let $G$ be the set of all $2^N$ configurations.
\begin{enumerate}
  \item Show that $G$ is an Abelian group isomorphic to $\Z_2 \times \dots \times \Z_2$ ($N$ factors). What are the identity and the inverse of $\sigma$?
  \item The bond between spins $i$ and $i + 1$ has the value $b_i = \sigma_i \sigma_{i+1}$. Show that $\varphi(\sigma) = (b_1, \dots, b_N)$ defines a homomorphism $\varphi\colon G \to G$, and find $\ker\varphi$. What does multiplication by an element of $\ker\varphi$ do to a configuration?
  \item A bond with $b_i = -1$ is a domain wall. Use the first isomorphism theorem to find $|\im\varphi|$. Show that every $b \in \im\varphi$ satisfies $b_1 b_2 \cdots b_N = 1$, and conclude that $\im\varphi$ consists exactly of the bond patterns with an even number of domain walls. Why must a ring have an even number of domain walls?
  \item The Ising energy is $E(\sigma) = -J\sum_{i=1}^{N}\sigma_i\sigma_{i+1}$ with $J > 0$. Show that $E$ is constant on each coset of $\ker\varphi$, and express $E$ in terms of the number $w$ of domain walls. How many configurations have exactly $w$ domain walls? For $N = 4$, list the possible energies with the number of configurations of each, and check that the numbers add up to 16.
  \item An open chain has no bond between spins $N$ and 1, so $\varphi_{\text{open}}(\sigma) = (\sigma_1\sigma_2, \dots, \sigma_{N-1}\sigma_N)$ takes values in $\Z_2^{\,N-1}$. Find its kernel and show that it is onto. Can an open chain have an odd number of domain walls?
\end{enumerate}
\showsolution{16}

%======================================================================
\section*{\large Part C. Linear algebra and representations (Lectures 7--8)}

\poolproblem{The vibrations of carbon dioxide}{Problem Set 3, Problem 1}
The carbon dioxide molecule is straight, O=C=O, with the carbon atom (mass $M$) between two oxygen atoms (mass $m$). Consider motion along the axis only: $x_1$, $x_2$ and $x_3$ are the displacements of the three atoms from equilibrium, and each bond acts as a spring with constant $k$:
\[
  m\ddot x_1 = -k(x_1 - x_2), \qquad M\ddot x_2 = -k(x_2 - x_1) - k(x_2 - x_3), \qquad m\ddot x_3 = -k(x_3 - x_2).
\]
\begin{enumerate}
  \item Write these equations as $\MM\ddot x = -Kx$, with $x = (x_1, x_2, x_3)^T$ and the mass matrix $\MM = \diag(m, M, m)$, and find the symmetric matrix $K$. Show that $x(t) = a\cos\omega t$, with a fixed vector $a \ne 0$, solves the equations exactly when $Ka = \omega^2\MM a$, that is, when $a$ is an eigenvector of $A = \MM^{-1}K$ with eigenvalue $\omega^2$. Show that $a^T K a = k\bigl[(a_1 - a_2)^2 + (a_2 - a_3)^2\bigr]$, and deduce that $\omega^2 \ge 0$ for every real mode $a$.
  \item $A$ is not symmetric when $M \ne m$. Show that $\langle a, b\rangle_\MM = a^T \MM b$ is an inner product on $\R^3$ for which $\langle a, Ab\rangle_\MM = \langle Aa, b\rangle_\MM$, and deduce that two modes with different frequencies are orthogonal for $\langle\cdot, \cdot\rangle_\MM$.
  \item The inversion through the carbon atom exchanges the two oxygen atoms and reverses the direction of the axis, so it maps a displacement $x$ to $Px = (-x_3, -x_2, -x_1)^T$. Write $P$ as a matrix, and show that $P^2 = I$ and that $P$ commutes with $\MM$ and $K$, and hence with $A$. Thus $\{I, P\}$ is a representation of $\Z_2$ on $\R^3$. Use its character to find how many times the trivial and the sign representations of $\Z_2$ appear in it. Find the subspaces of $\R^3$ on which $P = +1$ and $P = -1$, and show that $A$ maps each of them into itself.
  \item Find the three frequencies and modes: on the line where $P = +1$ directly, and on the plane where $P = -1$ by writing $A$ there as a $2 \times 2$ matrix in the basis $(1, 0, 1)^T$, $(0, 1, 0)^T$. Describe each mode in words. Why is one frequency zero? Check that the other two modes are orthogonal to it for $\langle\cdot, \cdot\rangle_\MM$, and explain what this means for the total momentum of the molecule.
  \item The oxygen atoms carry a small charge $-q$ each and the carbon atom $+2q$, so a displacement gives the molecule the electric dipole moment $d(x) = q(2x_2 - x_1 - x_3)$ along its axis. Show that $d(Px) = -d(x)$, and deduce that $d = 0$ at all times in a mode with $P = +1$. Only a vibration that makes the dipole oscillate absorbs infrared light: which of your modes does? This asymmetric stretch of CO$_2$ absorbs near a wavelength of $4.3\,\mu$m, and its bending vibrations (sideways, not in this model) near $15\,\mu$m, close to the peak of the infrared radiation emitted by the Earth. Why can a molecule of two identical atoms, like N$_2$ or O$_2$, the main gases of the air, not absorb infrared light in this way?
\end{enumerate}
\showsolution{17}

\poolproblem{Rotations and reflections of the plane: the group $O(2)$}{}
Let $O(2)$ be the set of real $2 \times 2$ matrices $O$ with $O^T O = I$. Let
\[
  R(\theta) = \begin{pmatrix} \cos\theta & -\sin\theta \\ \sin\theta & \cos\theta \end{pmatrix}, \qquad
  F = \begin{pmatrix} 1 & 0 \\ 0 & -1 \end{pmatrix}
\]
be the counterclockwise rotation by $\theta$ and the reflection in the $x$-axis.
\begin{enumerate}
  \item Show that $O(2)$ is a group under matrix multiplication, that $\det O = \pm 1$ for $O \in O(2)$, and that $\det\colon O(2) \to \{1, -1\}$ is a homomorphism. Its kernel $SO(2)$ is therefore a normal subgroup. What is $O(2)/SO(2)$?
  \item Show that every $O \in O(2)$ is either $R(\theta)$ or $R(\theta)F$ for some $\theta$. (Hint: the first column is a unit vector $(\cos\theta, \sin\theta)^T$, and the second is a unit vector orthogonal to it.) Which of the two has determinant 1?
  \item Show that
  \[
    R(\theta)F = \begin{pmatrix} \cos\theta & \sin\theta \\ \sin\theta & -\cos\theta \end{pmatrix},
  \]
  that $\bigl(R(\theta)F\bigr)^2 = I$, and that $\bigl(\cos\frac\theta2, \sin\frac\theta2\bigr)^T$ and $\bigl(-\sin\frac\theta2, \cos\frac\theta2\bigr)^T$ are eigenvectors with eigenvalues $+1$ and $-1$. (Use $\cos(a - b) = \cos a\cos b + \sin a \sin b$ and $\sin(a - b) = \sin a \cos b - \cos a \sin b$.) Conclude that $R(\theta)F$ is the reflection in the line at the angle $\theta/2$ to the $x$-axis.
  \item Show that $F R(\theta) F^{-1} = R(-\theta)$ and $R(\varphi)\,\bigl(R(\theta)F\bigr)\,R(\varphi)^{-1} = R(\theta + 2\varphi)F$, and interpret both results geometrically. Which elements of $O(2)$ are conjugate to $R(\theta)$, and which to $F$?
  \item Over $\C$, find the eigenvalues of $R(\theta)$ and show that $(1, -i)^T$ and $(1, i)^T$ are eigenvectors of $R(\theta)$ for every $\theta$. With
  \[
    S = \begin{pmatrix} 1 & 1 \\ -i & i \end{pmatrix},
  \]
  compute $S^{-1}R(\theta)S$ and $S^{-1}FS$. The rotations form a representation of $SO(2)$ on $\R^2$ with no invariant real line (apart from the rotations by $0$ and $\pi$). Is it irreducible over $\C$? Is the representation of all of $O(2)$ on $\C^2$ irreducible?
\end{enumerate}
\showsolution{18}

\poolproblem{Spin one-half: the Pauli matrices and $SU(2)$}{}
The Pauli matrices are
\[
  \sigma_x = \begin{pmatrix} 0 & 1 \\ 1 & 0 \end{pmatrix}, \qquad
  \sigma_y = \begin{pmatrix} 0 & -i \\ i & 0 \end{pmatrix}, \qquad
  \sigma_z = \begin{pmatrix} 1 & 0 \\ 0 & -1 \end{pmatrix}.
\]
\begin{enumerate}
  \item Show that each $\sigma_a$ is Hermitian and unitary, with $\sigma_a^2 = I$, $\tr\sigma_a = 0$ and $\det\sigma_a = -1$, and find its eigenvalues. Find normalized eigenvectors of $\sigma_x$ and of $\sigma_y$, and check that the two eigenvectors of each are orthogonal.
  \item Show that $\sigma_x\sigma_y = i\sigma_z$, $\sigma_y\sigma_z = i\sigma_x$ and $\sigma_z\sigma_x = i\sigma_y$, and that two different Pauli matrices anticommute: $\sigma_x\sigma_y = -\sigma_y\sigma_x$, and so on.
  \item Show that $U(2) = \{U : U^\dagger U = I\}$ is a group, that $|\det U| = 1$ for $U \in U(2)$, and that $SU(2) = \{U \in U(2) : \det U = 1\}$ is a subgroup. Show that every element of $SU(2)$ has the form
  \[
    U = \begin{pmatrix} a & -b^* \\ b & a^* \end{pmatrix}, \qquad a, b \in \C, \quad |a|^2 + |b|^2 = 1.
  \]
  (Hint: the columns of a unitary matrix are orthonormal.)
  \item For real $\theta$ let $U(\theta) = \cos\frac\theta2\, I - i\sin\frac\theta2\,\sigma_z$. Show that $U(\theta) = \diag\bigl(e^{-i\theta/2}, e^{i\theta/2}\bigr)$, that $U(\theta) \in SU(2)$, and that $\theta \mapsto U(\theta)$ is a homomorphism from $(\R, +)$ to $SU(2)$. Find its kernel. What is $U(2\pi)$?
  \item Show that
  \begin{gather*}
    U(\theta)\sigma_x U(\theta)^\dagger = \cos\theta\,\sigma_x + \sin\theta\,\sigma_y, \qquad
    U(\theta)\sigma_y U(\theta)^\dagger = -\sin\theta\,\sigma_x + \cos\theta\,\sigma_y, \\
    U(\theta)\sigma_z U(\theta)^\dagger = \sigma_z.
  \end{gather*}
  Thus, if a vector $(x, y, z)$ is written as the matrix $x\sigma_x + y\sigma_y + z\sigma_z$, conjugation by $U(\theta)$ rotates it by $\theta$ about the $z$-axis. Which two values of $U(\theta)$ give the identity rotation? (In Week 9 this becomes the two-to-one map $SU(2) \to SO(3)$: a spin one-half state changes sign under a rotation by $2\pi$.)
\end{enumerate}
\showsolution{19}

\poolproblem{Invariant complements: Galilean and Lorentz boosts}{}
\begin{enumerate}
  \item Let $D(e) = I$ and
  \[
    D(s) = \begin{pmatrix} 1 & 1 \\ 0 & -1 \end{pmatrix}.
  \]
  Show that this is a representation of $\Z_2 = \{e, s\}$ on $\R^2$, and find all invariant lines. Show that the line $W$ spanned by $e_1 = (1, 0)^T$ is invariant but its orthogonal complement $W^\perp$ is not, and find an invariant complement of $W$. Why does this not contradict the lecture result that $W^\perp$ is invariant?
  \item Compute $H = \sum_{g \in \Z_2} D(g)^T D(g)$ and check that $D(s)^T H D(s) = H$. Show that the complement of $W$ for the inner product $\langle u, v\rangle_H = u^T H v$, namely $\{u : e_1^T H u = 0\}$, is the invariant line you found in (a). Find an upper triangular $C$ with $C^T C = H$, and compute $C D(s) C^{-1}$. What kind of matrix is it?
  \item A Galilean boost with velocity $v$ acts on the space-time coordinates $(x, t)^T$ of an event as
  \[
    B(v) = \begin{pmatrix} 1 & v \\ 0 & 1 \end{pmatrix}.
  \]
  Show that $B(u)B(v) = B(u + v)$, so that $B$ is a representation of $(\R, +)$. Show that the line $t = 0$ is invariant, and that for $v \ne 0$ it is the only line invariant under $B(v)$. Is this representation a direct sum of two one-dimensional representations? What does the invariant line say about simultaneity?
  \item Show that no symmetric, positive-definite $H$ satisfies $B(v)^T H B(v) = H$ for every $v$. (Hint: compare the entries in the first row.) Why does the finite-group unitarity theorem not apply here?
  \item With $c = 1$, a Lorentz boost with rapidity $\eta$ acts on $(t, x)^T$ as
  \[
    L(\eta) = \begin{pmatrix} \cosh\eta & \sinh\eta \\ \sinh\eta & \cosh\eta \end{pmatrix}.
  \]
  Show that $L(\alpha)L(\beta) = L(\alpha + \beta)$, using $\cosh(\alpha + \beta) = \cosh\alpha\cosh\beta + \sinh\alpha\sinh\beta$ and $\sinh(\alpha + \beta) = \sinh\alpha\cosh\beta + \cosh\alpha\sinh\beta$. Find two invariant lines and the one-dimensional representations of $(\R, +)$ on them, and interpret the lines physically. Is $L$ a direct sum of one-dimensional representations? Are they unitary?
\end{enumerate}
\showsolution{20}

\poolproblem{The $\pi$ electrons of benzene}{}
In the simplest (H\"uckel) model of the $\pi$ electrons of benzene, an electron hops between the six carbon atoms of the ring, labelled $1, \dots, 6$ counterclockwise. Its state is a vector $c = (c_1, \dots, c_6)^T \in \C^6$, and its energies are the eigenvalues of
\[
  H = -t\,(T + T^{-1}), \qquad t > 0,
\]
where $T$ is the $6 \times 6$ permutation matrix with $Te_j = e_{j+1}$ (and $e_7 = e_1$): $T = P(r)$ for the rotation $r = (123456)$.
\begin{enumerate}
  \item Show that $T^6 = I$ and $T^{-1} = T^T$, and that $H$ is real, symmetric and commutes with $T$. Write $H$ as a matrix.
  \item Let $\omega = e^{i\pi/3}$. Show that the vector $v_k$ with components $(v_k)_j = \omega^{-kj}$, $j = 1, \dots, 6$, satisfies $T v_k = \omega^k v_k$ for $k = 0, 1, \dots, 5$. Explain why the six vectors $v_k$ form a basis of $\C^6$, and describe $\C^6$ as a direct sum of one-dimensional representations of $\Z_6 = \gen{r}$.
  \item Show that $Hv_k = E_k v_k$ with $E_k = -2t\cos(\pi k/3)$, and list the energies with their degeneracies. (Use $\cos 0 = 1$, $\cos 60^\circ = \tfrac12$, $\cos 120^\circ = -\tfrac12$ and $\cos 180^\circ = -1$.)
  \item Benzene has six $\pi$ electrons, and each level holds two (spin up and spin down). Fill the lowest levels and find the total $\pi$ energy. Compare with three separate double bonds, each described by two atoms with $H = -t\,\sigma_x$, $\sigma_x = \begin{pmatrix} 0 & 1 \\ 1 & 0 \end{pmatrix}$, and two electrons. How much energy does the delocalization over the ring gain?
  \item The levels $E_k$ and $E_{6-k}$ coincide. The reflection $s = (26)(35)$ sends atom $j$ to atom $2 - j$ (mod 6). Show that $P(s)$ commutes with $H$ and that $P(s)\,v_k = \omega^{-2k}\, v_{6-k}$. Explain why this forces $E_k = E_{6-k}$.
\end{enumerate}
\showsolution{21}

\poolproblem{The square: a two-dimensional representation of $D_4$}{}
The vertices $1, 2, 3, 4$ of a square sit at $(1, 0)$, $(0, 1)$, $(-1, 0)$ and $(0, -1)$, and
\[
  D(r) = \begin{pmatrix} 0 & -1 \\ 1 & 0 \end{pmatrix}, \qquad D(s) = \begin{pmatrix} 1 & 0 \\ 0 & -1 \end{pmatrix}
\]
are the counterclockwise quarter-turn and the reflection in the axis through vertex 1.
\begin{enumerate}
  \item Check that $D(r)^4 = D(s)^2 = I$ and $D(s)D(r)D(s) = D(r)^{-1}$, and explain why $D(r^k s^m) = D(r)^k D(s)^m$ defines a representation of $D_4$. Write all eight matrices. Is $D$ faithful? Which line is the mirror of $rs$?
  \item The conjugacy classes of $D_4$ are $\{e\}$, $\{r^2\}$, $\{r, r^3\}$, $\{s, r^2 s\}$ and $\{rs, r^3 s\}$. Compute the character of $D$ on each class, and check that it is constant on each class.
  \item Show that $D$ is irreducible over $\R$ and over $\C$. (Hint: an invariant line must be an eigenline of $D(s)$.)
  \item Restrict $D$ to the subgroup $\gen{r} \cong \Z_4$. Find $S$ with $S^{-1}D(r)S = \diag(i, -i)$, and say which of the one-dimensional representations $D_k(r^m) = i^{km}$ of $\Z_4$ appear. Is the restriction irreducible over $\C$? Over $\R$?
  \item Show that
  \[
    D'(r) = \begin{pmatrix} i & 0 \\ 0 & -i \end{pmatrix}, \qquad D'(s) = \begin{pmatrix} 0 & 1 \\ 1 & 0 \end{pmatrix}
  \]
  also defines a representation of $D_4$, and that it is equivalent to $D$: find $S$ with $D'(g) = S^{-1}D(g)S$ for every $g \in D_4$.
\end{enumerate}
\showsolution{22}

\poolproblem{Quadratic functions in square and triangular crystals}{}
A symmetry of the plane with matrix $M$ acts on functions by $(O_M f)(x, y) = f\bigl(M^{-1}(x, y)^T\bigr)$: the function is carried along with the plane. Let $V$ be the three-dimensional space of quadratic functions, with basis $x^2$, $xy$, $y^2$.
\begin{enumerate}
  \item Show that $O_{M_1} O_{M_2} = O_{M_1 M_2}$ and that $O_M$ maps $V$ into $V$. Thus every representation $g \mapsto M(g)$ of a group by $2 \times 2$ matrices gives a representation $g \mapsto O_{M(g)}$ on $V$.
  \item Square crystal: let $D_4$ act on the plane by the counterclockwise quarter-turn and the reflection in the $x$-axis,
  \[
    D(r) = \begin{pmatrix} 0 & -1 \\ 1 & 0 \end{pmatrix}, \qquad D(s) = \begin{pmatrix} 1 & 0 \\ 0 & -1 \end{pmatrix}.
  \]
  Find the $3 \times 3$ matrices of $O_{D(r)}$ and $O_{D(s)}$ in the basis $(x^2, xy, y^2)$. Show that the lines spanned by $x^2 + y^2$, $x^2 - y^2$ and $xy$ are invariant, and give the one-dimensional representation of $D_4$ on each by its values on $r$ and $s$.
  \item Compute the character of $V$ on the classes $\{e\}$, $\{r^2\}$, $\{r, r^3\}$, $\{s, r^2 s\}$, $\{rs, r^3 s\}$ of $D_4$, and check that it is the sum of the characters of the three one-dimensional representations in (b).
  \item Triangular crystal: let $D_3$ act by the plane representation $D_E(r) = R$, $D_E(s) = F$ of the Conventions. Show that $x^2 + y^2$ is invariant, that the plane spanned by $x^2 - y^2$ and $2xy$ is invariant, and that in the basis $(x^2 - y^2, 2xy)$ the matrices of $O_R$ and $O_F$ are $R^2$ and $F$. Show that this two-dimensional representation is irreducible and equivalent to $D_E$.
  \item Compute the character of $V$ for $D_3$ on the classes $\{e\}$, $\{r, r^2\}$, $\{s, rs, r^2 s\}$, and check that it is $\chi_1 + \chi_E$. (In the crystal-field theory of Week 7, this is why the orbitals shaped like $x^2 - y^2$ and $xy$ behave differently in square and in triangular crystals.)
\end{enumerate}
\showsolution{23}

%======================================================================
\section*{\large Part D. Schur's lemma, orthogonality and characters (Lectures 9--10)}

\poolproblem{Why polyacetylene is not a metal}{Problem Set 3, Problem 2}
Polyacetylene is a long chain of carbon atoms, each bonded to one hydrogen atom, with alternating double and single bonds between the carbons. In the simplest (H\"uckel) model of its $\pi$ electrons, take a ring of $N \ge 3$ cells, each containing two carbon atoms $A_n$ and $B_n$, with $n = 1, \dots, N$ counted modulo $N$ (so $A_{N+1} = A_1$). A state of an electron is a vector in $\C^{2N}$ with basis vectors $e_{A_n}$, $e_{B_n}$, and its energies are the eigenvalues of the Hamiltonian
\[
  He_{A_n} = -t_1 e_{B_n} - t_2 e_{B_{n-1}}, \qquad He_{B_n} = -t_1 e_{A_n} - t_2 e_{A_{n+1}},
\]
where $t_1 > 0$ is the hopping amplitude across the double bond $A_nB_n$ and $t_2 > 0$ across the single bond $B_nA_{n+1}$. The translation by one cell is $Te_{A_n} = e_{A_{n+1}}$, $Te_{B_n} = e_{B_{n+1}}$.
\begin{enumerate}
  \item Show that $H$ is real and symmetric, that $T^N = I$, and that $HT = TH$. Thus the powers of $T$ form a representation of $\Z_N$ on $\C^{2N}$, and $H$ commutes with every matrix in it.
  \item Use Schur's lemma to show that every irreducible representation of an Abelian group is one-dimensional. Show that $\Z_N = \{e, a, \dots, a^{N-1}\}$ has exactly $N$ inequivalent irreducible representations, $D(a^j) = \lambda^j$ with $\lambda^N = 1$.
  \item For $k = 2\pi l/N$ with $l = 0, 1, \dots, N - 1$, let
  \[
    u^A_k = \sum_{n=1}^{N} e^{ikn} e_{A_n}, \qquad u^B_k = \sum_{n=1}^{N} e^{ikn} e_{B_n}.
  \]
  Show that $Tu^A_k = e^{-ik}u^A_k$ and $Tu^B_k = e^{-ik}u^B_k$, and that the $2N$ vectors $u^A_k$, $u^B_k$ are orthogonal. (Hint: sum a geometric series.) Conclude that $\C^{2N}$ is the direct sum of the planes $V_k = \spn(u^A_k, u^B_k)$, and that each one-dimensional representation of $\Z_N$ appears in it exactly twice.
  \item Show that $H$ maps each $V_k$ into itself, with the matrix
  \[
    h(k) = -\begin{pmatrix} 0 & t_1 + t_2e^{-ik} \\ t_1 + t_2e^{ik} & 0 \end{pmatrix}
  \]
  in the basis $(u^A_k, u^B_k)$, and find its eigenvalues $E_\pm(k)$. The $2N \times 2N$ problem has become $N$ problems of size $2 \times 2$: this is Bloch's theorem. As a check, take $N = 3$ and $t_1 = t_2 = t$, a ring of six equal bonds like benzene, and list the six energies with their degeneracies.
  \item In the limit $N \to \infty$, which we take from now on, the allowed $k$ fill the interval $(-\pi, \pi]$ densely, and the energies form two bands. Find the range of each band. Every carbon atom gives one $\pi$ electron, and each level holds two electrons (spin up and spin down): show that the $2N$ electrons exactly fill the lower band. Show that for $t_1 = t_2$ the two bands touch, so an arbitrarily small energy lifts an electron into an empty level and the chain conducts like a metal, while for $t_1 \ne t_2$ there is an energy gap $2|t_1 - t_2|$ between them. Real polyacetylene has alternating bonds and starts to absorb light at about $1.5$\,eV. Taking this as the gap, what is $|t_1 - t_2|$? (The discovery that polyacetylene conducts well when doped with iodine was recognized by the 2000 Nobel Prize in Chemistry.)
\end{enumerate}
\showsolution{24}

\poolproblem{Degenerate modes of a square drum}{Problem Set 3, Problem 3}
The skin of a square drum covers $-1 \le x, y \le 1$ and is fixed along the edges. Its normal modes are vibrations $u(x, y)\cos\omega t$ with
\[
  u_{nm}(x, y) = f_n(x)f_m(y), \qquad f_n(x) = \sin\frac{n\pi(x + 1)}{2}, \qquad n, m = 1, 2, 3, \dots,
\]
and frequencies $\omega_{nm} = \frac{\pi c}{2}\sqrt{n^2 + m^2}$, where $c$ is the speed of waves on the skin. (A quantum particle in a square box has the same wavefunctions, with energies proportional to $n^2 + m^2$.) The symmetry group $D_4$ of the square acts on the plane by the quarter-turn and the reflection in the $x$-axis,
\[
  D(r) = \begin{pmatrix} 0 & -1 \\ 1 & 0 \end{pmatrix}, \qquad D(s) = \begin{pmatrix} 1 & 0 \\ 0 & -1 \end{pmatrix},
\]
and on functions by $(O_g u)(x, y) = u\bigl(D(g)^{-1}(x, y)^T\bigr)$; you may use that $O_{gh} = O_g O_h$. The character table of $D_4$ is
\[
  \begin{array}{c|ccccc}
    & \{e\} & \{r^2\} & \{r, r^3\} & \{s, r^2 s\} & \{rs, r^3 s\} \\ \hline
    1_{++} & 1 & 1 & 1 & 1 & 1 \\
    1_{+-} & 1 & 1 & 1 & -1 & -1 \\
    1_{-+} & 1 & 1 & -1 & 1 & -1 \\
    1_{--} & 1 & 1 & -1 & -1 & 1 \\
    E & 2 & -2 & 0 & 0 & 0
  \end{array}
\]
where $1_{\epsilon\epsilon'}$ is the one-dimensional representation with $r \mapsto \epsilon$ and $s \mapsto \epsilon'$, and $E$ is the representation $D$ itself.
\begin{enumerate}
  \item Show that $f_n(-x) = (-1)^{n+1}f_n(x)$. Show that
  \[
    O_r u_{nm} = (-1)^{m+1}u_{mn}, \qquad O_s u_{nm} = (-1)^{m+1}u_{nm},
  \]
  and find $O_{r^2}u_{nm}$ and $O_{rs}u_{nm}$. Deduce that every $O_g$ maps the modes of each frequency to modes of the same frequency, so that these modes span a representation of $D_4$.
  \item For $n \ne m$, the modes $u_{nm}$ and $u_{mn}$ have the same frequency. Find the character of the representation of $D_4$ on $W_{nm} = \spn(u_{nm}, u_{mn})$ on the five classes. Show that $W_{nm}$ is irreducible, and equivalent to $E$, exactly when $n + m$ is odd.
  \item Let $n \ne m$ with $n + m$ even. Decompose $W_{nm}$ with characters, find its two invariant lines, and say which one-dimensional representation acts on each. Treat the cases of $n$ and $m$ both odd and both even separately.
  \item A small change of the drum that keeps all the symmetries of the square, for example a small weight glued to its centre, shifts the frequencies. For the modes of one frequency, the first-order shifts of their frequencies are the eigenvalues of a real symmetric matrix $V$ that acts on their span and commutes with every $O_g$, and its eigenvectors are the new modes. Use Schur's lemma to show that $u_{12}$ and $u_{21}$ keep a common frequency, while $u_{13}$ and $u_{31}$ may split, and that the new modes are then $u_{13} \pm u_{31}$. Check this for the weight at the centre: a mode with a node at the centre is not affected, and a mode that moves there is slowed down. (Use $\sin\frac{\pi}{2} = 1$, $\sin\pi = 0$ and $\sin\frac{3\pi}{2} = -1$.) The degeneracy of $u_{13}$ and $u_{31}$ is called accidental: the symmetry of the square does not require it.
  \item The smallest value of $n^2 + m^2$ shared by three modes is $50 = 1^2 + 7^2 = 5^2 + 5^2$. Without computing, explain why the symmetry of the square cannot be the reason for this threefold degeneracy. Then find the character of the representation on the span of $u_{17}$, $u_{71}$, $u_{55}$, and decompose it.
\end{enumerate}
\showsolution{25}

\poolproblem{Three spins on a triangle}{Problem Set 3, Problem 4}
Three spin-one-half particles, for example the unpaired electrons of three magnetic ions at the corners of a triangle, sit at the sites $1, 2, 3$. Each spin is up ($\up$) or down ($\dn$), so a state is a vector in $\C^8$ with the basis $\ket{s_1s_2s_3}$, $s_i \in \{\up, \dn\}$. A permutation $p \in S_3$ moves the spin at site $i$ to site $p(i)$:
\[
  P(p)\ket{s_1s_2s_3} = \ket{s'_1s'_2s'_3}, \qquad s'_{p(i)} = s_i.
\]
For example, $P\bigl((12)\bigr)\ket{\up\dn\dn} = \ket{\dn\up\dn}$.
\begin{enumerate}
  \item Show that $P(pq) = P(p)P(q)$, and that $\chi_P(p)$ is the number of basis states left unchanged by $P(p)$. Find $\chi_P$ on the three classes of $S_3$, and decompose $P$ into irreducible representations.
  \item Permutations do not change the number of up spins, so the subspaces with 3, 2, 1 and 0 up spins (of dimensions 1, 3, 3 and 1) are invariant. Decompose each of them, and check that the results add up to (a). In the subspace with two up spins, find the line on which $S_3$ acts trivially and an invariant complement on which it acts by $E$.
  \item Three electrons in the same orbital, such as the 1s orbital of an atom, have a wavefunction that is symmetric in their positions, so the Pauli principle requires their spin state $\psi$ to be antisymmetric: $P(p)\psi = \sgn(p)\psi$ for every $p \in S_3$. Show that no such $\psi \ne 0$ exists. The third electron of a lithium atom therefore cannot occupy the 1s orbital; the lowest orbital still available to it is 2s.
  \item Let $K = P\bigl((12)\bigr) + P\bigl((23)\bigr) + P\bigl((13)\bigr)$. Show that $K$ commutes with every $P(p)$. More generally, for any representation $D$ of $S_3$ show that $K_D = D\bigl((12)\bigr) + D\bigl((23)\bigr) + D\bigl((13)\bigr)$ commutes with every $D(p)$, and use Schur's lemma to show that for irreducible $D$ of dimension $d$, $K_D = \lambda I$ with $\lambda = 3\chi_D(\tau)/d$, where $\tau$ is a transposition. Find $\lambda$ for $1$, $\sgn$ and $E$.
  \item The Heisenberg model of the three ions has the Hamiltonian $H = J(\vec S_1\cdot\vec S_2 + \vec S_2\cdot\vec S_3 + \vec S_3\cdot\vec S_1)$, where $\vec S_i$ is the spin of the particle at site $i$. For spin one-half, Dirac's identity $\vec S_i\cdot\vec S_j = \tfrac12 P\bigl((ij)\bigr) - \tfrac14$ (with $\hbar = 1$) turns it into
  \[
    H = \frac{J}{2}K - \frac{3J}{4},
  \]
  and you may take this as the definition of $H$. Without diagonalizing an $8 \times 8$ matrix, find the energies and their degeneracies. For $J > 0$ every pair of neighbours prefers opposite spins, which cannot happen for all three pairs of a triangle: the magnet is frustrated. Show that its lowest level is fourfold degenerate. A magnetic field along the $z$-axis adds $-hS_z$ to $H$, where $S_z$ is half the number of up spins minus half the number of down spins. Show that for $0 < h < \tfrac32 J$ the lowest level is still twofold degenerate, and explain with Schur's lemma why the field cannot remove this degeneracy.
\end{enumerate}
\showsolution{26}

\poolproblem{The vibrations of ammonia}{}
The ammonia molecule NH$_3$ is a pyramid. The nitrogen atom sits on the $z$-axis, and the hydrogen atoms $1, 2, 3$ at the corners of an equilateral triangle below it, at the angles $0^\circ$, $120^\circ$ and $240^\circ$ from the $x$-axis. Its symmetry group $G \cong S_3$ consists of the rotations by $0$ and $\pm 120^\circ$ about the $z$-axis and the reflections in the three vertical planes that contain the $z$-axis and one hydrogen atom. With $r = (123)$, the rotation by $120^\circ$, and $s = (23)$, the reflection in the $xz$-plane, $G$ acts on $\R^3$ by the vector representation
\[
  D_V(r) = \begin{pmatrix} R & 0 \\ 0 & 1 \end{pmatrix}, \qquad D_V(s) = \begin{pmatrix} F & 0 \\ 0 & 1 \end{pmatrix},
\]
with $R$ and $F$ as in the Conventions. Each of the four atoms $a = \text{N}, 1, 2, 3$ can be displaced by a vector $u_a \in \R^3$, so the displacements form a 12-dimensional real space. An element $g \in G$ carries atom $a$ to atom $g(a)$ and its displacement $u_a$ to $D_V(g)u_a$, which defines the representation $\Gamma$ on this space:
\[
  \bigl(\Gamma(g)u\bigr)_{g(a)} = D_V(g)\,u_a.
\]
\begin{enumerate}
  \item Find $\chi_V$ and show that $D_V \simeq 1 \oplus E$. Show that
  \[
    \chi_\Gamma(g) = \bigl(\text{number of atoms with } g(a) = a\bigr)\times\chi_V(g).
  \]
  (Hint: write $\Gamma(g)$ as a $4 \times 4$ array of $3 \times 3$ blocks.) Find $\chi_\Gamma$ on the three classes, and decompose $\Gamma$.
  \item The displacements with $u_a = b$ for every atom, $b \in \R^3$, are the translations of the molecule. Show that they form an invariant subspace on which $G$ acts by $D_V$. The small rotations of the molecule about its centre of mass form another three-dimensional invariant subspace. Show that its character is $(3, 0, -1)$ on the classes $\{e\}$, $\{r, r^2\}$, $\{s, rs, r^2 s\}$. (Hint: a reflection reverses the sense of a rotation about an axis that lies in its mirror, and keeps the sense of a rotation about an axis perpendicular to the mirror.) Decompose both subspaces.
  \item The vibrations are the displacements $u$ that are orthogonal to all translations and rotations for the mass-weighted inner product $\langle u, w\rangle = \sum_a m_a\, u_a \cdot w_a$; they form a six-dimensional invariant complement of the subspace in (b). Show that $G$ acts on the vibrations by $1 \oplus 1 \oplus E \oplus E$. The squared frequencies are the eigenvalues of a matrix that commutes with every $\Gamma(g)$ on this space. Explain why the eigenspace of each frequency is an invariant subspace, and deduce that a frequency whose eigenspace contains a copy of $E$ is at least twofold degenerate. If every eigenspace is irreducible (no accidental degeneracy), how many different frequencies does ammonia have, and how many of them are twofold degenerate? Compare with the frequencies measured in the infrared spectrum of ammonia: about $950$, $1627$, $3337$ and $3444\ \text{cm}^{-1}$.
  \item Find all displacements left unchanged by every $\Gamma(g)$, and check that they form a space of the dimension predicted by (a). (Hint: the nitrogen atom can only move along the $z$-axis, and each hydrogen atom moves by the same amount $\beta$ away from the axis and by the same amount $\gamma$ along it.) Which of them is a translation? Describe the umbrella motion, the one in which the hydrogen atoms move only along the $z$-axis and the centre of mass stays at rest. Continued far enough, this motion turns the pyramid inside out; quantum tunnelling between the two pyramids gives the 24\,GHz radiation of the ammonia maser, the forerunner of the laser.
\end{enumerate}
\showsolution{27}

\poolproblem{What commutes with a representation?}{}
Use the plane representation of $D_4$ by the quarter-turn and the reflection in the $x$-axis:
\[
  D(r) = \begin{pmatrix} 0 & -1 \\ 1 & 0 \end{pmatrix}, \qquad D(s) = \begin{pmatrix} 1 & 0 \\ 0 & -1 \end{pmatrix}.
\]
\begin{enumerate}
  \item Find all complex $2 \times 2$ matrices $A$ with $AD(r) = D(r)A$ and $AD(s) = D(s)A$. Which form of Schur's lemma does your answer illustrate, and what does it tell you about $D$?
  \item Find all complex $A$ that commute with $D(r)$ alone. Over $\R$, the representation of the subgroup $\gen{r} \cong \Z_4$ on $\R^2$ has no invariant line, yet $A$ need not be a multiple of $I$. Why does this not contradict Schur's lemma? Over $\C$, find the invariant lines of $\gen{r}$ and explain your answer with Schur's lemma.
  \item The four one-dimensional representations of $D_4$ have $c(r) = \pm 1$ and $c(s) = \pm 1$. For each of them, find all $2 \times 1$ matrices $A$ with $D(g)A = A\,c(g)$ for every $g \in D_4$. Explain the result with the intertwiner form of Schur's lemma.
  \item Let $z$ be an element that commutes with every element of a finite group $G$, and let $D$ be an irreducible complex representation of $G$. Show that $D(z) = \lambda I$ with $\lambda^{\ord(z)} = 1$. In $D_4$, $r^2$ commutes with every element: find $D(r^2)$ in the plane representation and $c(r^2)$ in the four one-dimensional ones, and check the claim.
  \item A Hamiltonian $H$ for two states of an electron in a square crystal, such as orbitals shaped like $x$ and $y$, acts on $\C^2$ and satisfies $HD(g) = D(g)H$ for every $g \in D_4$. Show that the two states have the same energy. If the crystal is stretched along the $x$-axis, its symmetry group becomes $\{e, r^2, s, r^2 s\}$. Find all $H$ that commute with $D(g)$ for these four $g$. Must the two energies still be equal?
\end{enumerate}
\showsolution{28}

\poolproblem{Symmetry and isotropy: Neumann's principle}{}
In a conducting crystal, the current density $J$ and the electric field $E$ are related by $J = \sigma E$, where the conductivity $\sigma$ is a real matrix. If a rotation or reflection $D(g)$ maps the crystal to itself, turning the field by $D(g)$ turns the current by $D(g)$, so $\sigma D(g) E = D(g) \sigma E$ for every $E$: the matrix $\sigma$ commutes with every $D(g)$.
\begin{enumerate}
  \item Square lattice in a plane: the symmetry group is $D_4$, acting by the quarter-turn and the reflection
  \[
    D(r) = \begin{pmatrix} 0 & -1 \\ 1 & 0 \end{pmatrix}, \qquad D(s) = \begin{pmatrix} 1 & 0 \\ 0 & -1 \end{pmatrix}.
  \]
  Using the classes $\{e\}$, $\{r^2\}$, $\{r, r^3\}$, $\{s, r^2 s\}$, $\{rs, r^3 s\}$, show with characters that this representation is irreducible over $\C$. Use Schur's lemma to show that $\sigma = \sigma_0 I$ with real $\sigma_0$: the current is parallel to the field in every direction, although the crystal itself is not isotropic.
  \item A triangular or hexagonal lattice, such as graphene, has the symmetries of $D_3$ acting by the plane representation $D_E$ of the Conventions (and more). Show that $D_E$ is irreducible and conclude again that $\sigma$ is a multiple of $I$.
  \item A rectangular lattice has the symmetry group $\{I, -I, F, -F\}$ with $F = \diag(1, -1)$. Find all $\sigma$ that commute with these matrices. Why does Schur's lemma not force isotropy here?
  \item For the square lattice, compute the group average
  \[
    \mathcal{A}(X) = \frac18 \sum_{g \in D_4} D(g)\, X\, D(g)^T, \qquad X = \begin{pmatrix} a & b \\ c & d \end{pmatrix},
  \]
  directly from the eight matrices, and check that $\mathcal{A}(X) = \frac{\tr X}{2} I$, as Lecture 9 predicts.
  \item Cubic crystals: the 24 rotations of a cube act on $\R^3$ by $3 \times 3$ rotation matrices, and a rotation by $\theta$ has trace $1 + 2\cos\theta$. The classes of this group are: the identity; the 6 quarter-turns about the three face axes; the 3 half-turns about these axes; the 8 turns by $\pm 120^\circ$ about the four body diagonals; the 6 half-turns about the six axes through midpoints of opposite edges. Using $\cos 90^\circ = 0$, $\cos 120^\circ = -\tfrac12$ and $\cos 180^\circ = -1$, show that this three-dimensional representation is irreducible. What does Schur's lemma then say about the conductivity of copper or rock salt, and about the moment of inertia tensor of a uniform cube about its centre?
\end{enumerate}
\showsolution{29}

\poolproblem{Orthogonality of matrix entries for $S_3$}{}
Order the elements of $S_3$ as $e, r, r^2, s, rs, r^2 s$, with $r = (123)$ and $s = (23)$. Its irreducible representations are the trivial representation, the sign representation and the plane representation $D_E$, with $D_E(r) = R$ and $D_E(s) = F$ as in the Conventions.
\begin{enumerate}
  \item Compute the six matrices $D_E(g)$ and check that each is orthogonal.
  \item Use your table to evaluate
  \begin{gather*}
    \frac16 \sum_g D_E(g)_{11} D_E(g)_{22}, \qquad \frac16 \sum_g D_E(g)_{12}^{\,2}, \\
    \frac16 \sum_g D_E(g)_{11} D_E(g)_{21}, \qquad \frac16 \sum_g \sgn(g)\, D_E(g)_{11},
  \end{gather*}
  and check each against the Great Orthogonality Theorem. Which Kronecker delta makes each vanishing sum vanish?
  \item A function on $S_3$ is a list of six complex numbers. Show that the six functions $1$, $\sgn$ and $\sqrt2\, D_E(g)_{ij}$ ($i, j = 1, 2$) form an orthonormal basis of the space of functions on $S_3$ with the inner product $\langle f, h\rangle_G$. Why does it matter that $1^2 + 1^2 + 2^2 = 6$?
  \item Expand the function $\delta_e$, with $\delta_e(e) = 1$ and $\delta_e(g) = 0$ for $g \ne e$, in this basis, and show that
  \[
    \delta_e(g) = \tfrac16\bigl(\chi_1(g) + \chi_{\sgn}(g) + 2\chi_E(g)\bigr).
  \]
  How is this related to the character of the regular representation?
  \item Use the Great Orthogonality Theorem to show that an irreducible representation of a group of order $N$ has dimension at most $\sqrt N$. Can a group of order 6 have an irreducible representation of dimension 3? Can $S_3$ have a second, inequivalent, two-dimensional irreducible representation?
\end{enumerate}
\showsolution{30}

\poolproblem{Counting isomers with characters}{}
A finite group $G$ permutes a finite set $X$; write $gx$ for the image of $x \in X$. The permutation representation acts on the space with one basis vector $e_x$ for each $x \in X$ by $P(g)e_x = e_{gx}$. The orbit of $x$ is $Gx = \{gx : g \in G\}$.
\begin{enumerate}
  \item Show that $\chi_P(g)$ is the number of points $x \in X$ with $gx = x$.
  \item Show that $v = \sum_x v_x e_x$ satisfies $P(g)v = v$ for every $g \in G$ exactly when $v_x$ is constant on each orbit. Deduce that these fixed vectors form a space whose dimension is the number of orbits.
  \item Show that this dimension is the multiplicity $m_1$ of the trivial representation in $P$, and deduce Burnside's counting formula:
  \[
    \text{number of orbits} = \frac{1}{|G|}\sum_{g \in G} \bigl(\text{number of points fixed by } g\bigr).
  \]
  \item In dichlorobenzene, two of the six hydrogen atoms of benzene are replaced by chlorine atoms. Two substitution patterns give the same compound if a symmetry of the hexagon, an element of $D_6$, maps one to the other. Let $X$ be the set of the 15 two-element subsets of $\{1, \dots, 6\}$. For each kind of element of $D_6$ (the identity; the rotations by $\pm 60^\circ$, by $\pm 120^\circ$ and by $180^\circ$; the three reflections whose mirrors pass through two atoms; the three whose mirrors pass through the midpoints of two bonds), count the subsets it fixes, and use (c) to find the number of different dichlorobenzenes. Draw one representative of each.
  \item Repeat (d) for trichlorobenzene, with three chlorine atoms.
\end{enumerate}
\showsolution{31}

\poolproblem{Spin one-half and the quaternion group}{}
With the Pauli matrices
\(
  \sigma_x = \bigl(\begin{smallmatrix} 0 & 1 \\ 1 & 0 \end{smallmatrix}\bigr),\
  \sigma_y = \bigl(\begin{smallmatrix} 0 & -i \\ i & 0 \end{smallmatrix}\bigr),\
  \sigma_z = \bigl(\begin{smallmatrix} 1 & 0 \\ 0 & -1 \end{smallmatrix}\bigr),
\)
let $A = -i\sigma_x$, $B = -i\sigma_y$ and $C = -i\sigma_z$: the $SU(2)$ matrices of half-turns of a spin one-half about the $x$-, $y$- and $z$-axes. Let $Q = \{\pm I, \pm A, \pm B, \pm C\}$.
\begin{enumerate}
  \item Show that $A^2 = B^2 = C^2 = -I$ and $AB = C$, $BC = A$, $CA = B$, $BA = -C$. Deduce that $Q$ is a non-Abelian group of order 8, and find the order of each element.
  \item Show that the conjugacy classes of $Q$ are $\{I\}$, $\{-I\}$, $\{\pm A\}$, $\{\pm B\}$ and $\{\pm C\}$.
  \item The defining representation $D(g) = g$ of $Q$ on $\C^2$ has the character $(2, -2, 0, 0, 0)$ on these classes. Show that it is irreducible.
  \item How many inequivalent irreducible representations does $Q$ have, and of which dimensions? Find the one-dimensional ones (hint: for $c\colon Q \to \C^\times$, show that $c(-I) = 1$ and $c(A), c(B) = \pm 1$), and write the character table of $Q$.
  \item $D_4$ also has classes of sizes $1, 1, 2, 2, 2$: $\{e\}$, $\{r^2\}$, $\{r, r^3\}$, $\{s, r^2 s\}$ and $\{rs, r^3 s\}$. Its irreducible representations are the four one-dimensional ones with $c(r) = \pm 1$, $c(s) = \pm 1$ and the plane representation of the square, with the character $(2, -2, 0, 0, 0)$. Write its character table and show that, with a suitable matching of the classes, it coincides with that of $Q$. Show nevertheless that $Q$ and $D_4$ are not isomorphic. What does this say about how much a character table tells us about a group?
\end{enumerate}
\showsolution{32}

\poolproblem{Rotations of a tetrahedron: the character table of $A_4$}{}
The 12 rotations that map a regular tetrahedron to itself, for example the four hydrogen atoms of a methane molecule, form a group $T \cong A_4$. Its conjugacy classes (Week 3) are $C_1 = \{e\}$; $C_2$, the 3 half-turns about the axes through midpoints of opposite edges; $C_3$, the 4 rotations by $120^\circ$ about the axes through a vertex and the centre of the opposite face, counterclockwise when seen from the vertex; and $C_4$, the 4 rotations by $240^\circ$ about the same axes. The identity and the half-turns form a normal subgroup $V \cong \Z_2 \times \Z_2$, and $T/V \cong \Z_3$. Let $\omega = e^{2\pi i/3}$, and remember that the inner product conjugates its first argument.
\begin{enumerate}
  \item How many inequivalent irreducible representations does $T$ have, and of which dimensions?
  \item Use $T/V \cong \Z_3$ to find three one-dimensional representations, and give their characters on $C_1, \dots, C_4$. (Hint: follow $g \mapsto gV$ by a representation of $\Z_3$.)
  \item The rotation matrices of $T$ act on $\R^3$, and a rotation by $\theta$ has trace $1 + 2\cos\theta$. Find the character of this representation (use $\cos 120^\circ = \cos 240^\circ = -\tfrac12$ and $\cos 180^\circ = -1$) and show that it is irreducible. Complete the character table of $T$ and check that its rows are orthonormal.
  \item $T$ permutes the four vertices of the tetrahedron. Find the character of this permutation representation and decompose it. In methane, the stretchings of the four C--H bonds combine into one symmetric ``breathing'' vibration and three vibrations of equal frequency. Explain the equal frequencies with Schur's lemma.
  \item $T$ also permutes the six edges of the tetrahedron. Find the character of this permutation representation and decompose it, and check the dimension.
\end{enumerate}
\showsolution{33}

\poolproblem{Rotations of a cube: decomposing permutation representations}{}
The 24 rotations of a cube form the group $O \cong S_4$ (Problem 14). Its classes, with their sizes, are the identity (1); the quarter-turns by $\pm 90^\circ$ about the three face axes (6); the half-turns about the face axes (3); the turns by $\pm 120^\circ$ about the four body diagonals (8); and the half-turns about the six axes through midpoints of opposite edges (6). With the names used in chemistry, the character table of $O$ is
\[
  \begin{array}{c|ccccc}
    & e & 6\ \text{quarter-turns} & 3\ \text{face half-turns} & 8\ \text{turns by } 120^\circ & 6\ \text{edge half-turns} \\ \hline
    A_1 & 1 & 1 & 1 & 1 & 1 \\
    A_2 & 1 & -1 & 1 & 1 & -1 \\
    E & 2 & 0 & 2 & -1 & 0 \\
    T_1 & 3 & 1 & -1 & 0 & -1 \\
    T_2 & 3 & -1 & -1 & 0 & 1
  \end{array}
\]
\begin{enumerate}
  \item Check that the squares of the dimensions add up to 24, that there are as many irreducible representations as classes, and that the rows $T_1$ and $T_2$ are orthonormal.
  \item The $3 \times 3$ rotation matrices form the vector representation. A rotation by $\theta$ has trace $1 + 2\cos\theta$. Using $\cos 90^\circ = 0$, $\cos 120^\circ = -\tfrac12$ and $\cos 180^\circ = -1$, find its character and identify it in the table.
  \item The rotations permute the 6 faces of the cube. Find the character of this permutation representation and decompose it. (These are the combinations of six orbitals on the neighbours of an atom at the centre of an octahedron, as in crystal-field theory.)
  \item Do the same for the 8 vertices of the cube.
  \item The rotations also permute the 4 body diagonals, which gives the isomorphism $O \cong S_4$. Decompose this permutation representation, and show that its three-dimensional piece is not equivalent to the vector representation, although both are irreducible and three-dimensional.
\end{enumerate}
\showsolution{34}

\poolproblem{A square molecule: degeneracy from Schur's lemma}{}
Four identical atoms sit at the vertices $1, 2, 3, 4$ of a square, labelled counterclockwise. An electron hops between neighbouring atoms with amplitude $t$ and across the diagonals with amplitude $t'$, so its Hamiltonian on $\C^4$ is
\[
  H = -t\,\bigl(P(r) + P(r)^{-1}\bigr) - t'\,P(r^2),
\]
where $P$ is the permutation representation of $D_4$ on the vertices, $P(g)e_j = e_{g(j)}$, with $r = (1234)$ and $s = (24)$.
\begin{enumerate}
  \item Show that $H$ commutes with $P(g)$ for every $g \in D_4$.
  \item $D_4$ has the five conjugacy classes $\{e\}$, $\{r^2\}$, $\{r, r^3\}$, $\{s, r^2 s\}$ and $\{rs, r^3 s\}$. Show that $D_4$ has exactly five inequivalent irreducible representations: four one-dimensional and one two-dimensional.
  \item Find the character of $P$, and compute $\langle\chi_P, \chi_P\rangle_G$ and the multiplicity of the trivial representation. Without using the rest of the character table, deduce that $P \simeq 1 \oplus A \oplus E$, where $A$ is a nontrivial one-dimensional representation and $E$ is the two-dimensional one, each appearing once.
  \item Use Schur's lemma to show that $H$ has at most three different eigenvalues, one of them at least twofold degenerate, whatever the values of $t$ and $t'$. Find the energies by applying $H$ to $(1, 1, 1, 1)^T$, $(1, -1, 1, -1)^T$, $(1, 0, -1, 0)^T$ and $(0, 1, 0, -1)^T$. Give $A$ by its values on $r$ and $s$.
  \item The square is stretched along the $x$-axis, through vertices 1 and 3, into a rhombus with four equal edges, and the hopping amplitudes across the two diagonals become different, $t'_{13} \ne t'_{24}$. Find the symmetry group of the rhombus, and the energies of $(1, 0, -1, 0)^T$ and $(0, 1, 0, -1)^T$ now. Explain with characters why the degeneracy can now be lifted.
\end{enumerate}
\showsolution{35}

\end{document}
